\documentclass{article}
\usepackage[T1]{fontenc}
\usepackage{lmodern}
\usepackage{graphicx}
\usepackage{amsmath,amssymb}
\usepackage[a4paper,margin=2cm]{geometry}
\usepackage[absolute,overlay]{textpos}
\setlength{\TPHorizModule}{1mm}
\setlength{\TPVertModule}{1mm}
\usepackage[indonesian]{babel}
\usepackage{hyperref}
\setlength{\emergencystretch}{2em}
% unit: Halaman 4

\begin{document}

% === Halaman 4 (llm) ===

\section*{One-Time Pad (OTP)}

\begin{itemize}
    \item Satu-satunya algoritma kriptografi sempurna aman (\textit{perfect secrecy}) sehingga tidak dapat dipecahkan adalah \textit{one-time pad (OTP)}.
    \item OTP ditemukan oleh Gilbert Vernam pada tahun 1917 dan kemudian secara matematis dibuktikan memiliki \textit{perfect secrecy} oleh Claude Shannon tahun 1949.
    \item OTP mengatasi kelemahan pada pada \textit{Vigenere Cipher}. \textit{Vigenere Cipher} mengulang penggunaan kunci secara periodik $\rightarrow$ mudah ditemukan dengan metode Kasiski.
    \item Pada OTP, panjang kunci = panjang plainteks, seluruh karakternya acak sejati
\end{itemize}

\vspace{5mm}

\begin{flushleft}
\textcolor{red}{\textbf{Plainteks:}} \texttt{otpadalahcipheryangtidakbisadipecahkan} \\
\textcolor{red}{\textbf{Kunci:}} \texttt{trjkdndkdwerylgrgdkopcegyhbdwjbtrfhgvk}
\end{flushleft}

\end{document}
